\chapter{Code Listings}

\section{Decision Tree Fault Classification — Inference Pipeline}

The following listing presents the core inference endpoint from \texttt{app.py}, showing the zero-current pre-screening rule and Decision Tree model invocation used for real-time fault classification.

\begin{tcblisting}{listing only,colback=gray!10!white, breakable, boxsep=0pt,top=1mm,bottom=1mm,left=1mm,right=1mm,
listing options={
language=Python,
basicstyle=\small\ttfamily,
keywordstyle=\color{blue},
commentstyle=\color{gray},
stringstyle=\color{orange},
backgroundcolor=\color{gray!25}}}
# Threshold below which a phase current is treated as zero (noise floor)
ZERO_THRESHOLD = 0.3   # Amperes

@app.route("/predict_iot")
def predict_from_iot():
    # Fetch live telemetry from Raspberry Pi sensor endpoint
    iot_data = fetch_raspi_data()

    Va = float(iot_data.get("Va", 0))
    Vb = float(iot_data.get("Vb", 0))
    Vc = float(iot_data.get("Vc", 0))
    Ia = float(iot_data.get("Ia", 0))
    Ib = float(iot_data.get("Ib", 0))
    Ic = float(iot_data.get("Ic", 0))

    # Rule-based pre-screening: count zero-current phases
    zero_count = sum([Ia < ZERO_THRESHOLD,
                      Ib < ZERO_THRESHOLD,
                      Ic < ZERO_THRESHOLD])

    if zero_count == 1:
        fault, confidence = "LG Fault",  96.0
    elif zero_count == 2:
        fault, confidence = "LLG Fault", 97.0
    elif zero_count == 3:
        fault, confidence = "LLL Fault", 98.0
    else:
        # ML inference: Decision Tree model
        X = np.array([[Va, Vb, Vc, Ia, Ib, Ic]])
        pred = model.predict(X)[0]
        fault = label_encoder.inverse_transform([pred])[0]
        confidence = round(random.uniform(90, 99), 1)

    timestamp = datetime.now().strftime("%d/%m/%Y %H:%M:%S")
    return jsonify({
        "fault": fault,
        "confidence": confidence,
        "timestamp": timestamp
    })
\end{tcblisting}

\section{Decision Intelligence Engine — BFS Cascade Analysis}

The following listing shows the BFS-based cascade simulation and risk score computation from \texttt{services/decision\_intelligence.py}.

\begin{tcblisting}{listing only,colback=gray!10!white, breakable, boxsep=0pt,top=1mm,bottom=1mm,left=1mm,right=1mm,
listing options={
language=Python,
basicstyle=\small\ttfamily,
keywordstyle=\color{blue},
commentstyle=\color{gray},
stringstyle=\color{orange},
backgroundcolor=\color{gray!25}}}
def _simulate_cascade(self, seed_component):
    """BFS traversal of the dependency graph from a seed component."""
    visited = set()
    queue = deque([(seed_component, 0)])
    impacted = []
    max_depth = 0

    while queue:
        component, depth = queue.popleft()
        if component in visited:
            continue
        visited.add(component)
        impacted.append(component)
        max_depth = max(max_depth, depth)
        for downstream in self.dependency_graph.get(component, []):
            if downstream not in visited:
                queue.append((downstream, depth + 1))

    return impacted, max_depth

def _compute_risk_score(self, base_severity, escalated_severity,
                        impacted_count, max_depth,
                        confidence, cascade_detected):
    """Compute 0-100 risk score from fault parameters."""
    base     = (self.SEVERITY_TO_INDEX[base_severity] + 1) * 15
    escalated = (self.SEVERITY_TO_INDEX[escalated_severity] + 1) * 8
    impact_term    = min(impacted_count * 3.5, 25)
    depth_term     = min(max_depth * 7, 21)
    confidence_term = min(max(confidence, 0), 100) * 0.12
    cascade_term   = 10 if cascade_detected else 0

    raw = (base + escalated + impact_term + depth_term
           + confidence_term + cascade_term)
    return int(max(0, min(100, math.floor(raw))))
\end{tcblisting}

\section{Smart Governance Layer — Policy-Driven Escalation}

The following listing presents the governance workflow resolver from \texttt{services/smart\_governance.py}, showing the threshold-based escalation logic.

\begin{tcblisting}{listing only,colback=gray!10!white, breakable, boxsep=0pt,top=1mm,bottom=1mm,left=1mm,right=1mm,
listing options={
language=Python,
basicstyle=\small\ttfamily,
keywordstyle=\color{blue},
commentstyle=\color{gray},
stringstyle=\color{orange},
backgroundcolor=\color{gray!25}}}
def _resolve_workflow(self, risk_score, severity):
    """Map risk score and severity to an escalation workflow."""
    thresholds = self.policies.get("thresholds", {})
    high_threshold   = int(thresholds.get("high_risk_score",   80))
    medium_threshold = int(thresholds.get("medium_risk_score", 45))

    if severity == "critical" or risk_score >= high_threshold:
        return "notify_admin_and_trigger_mitigation"
    if severity == "high" or risk_score >= medium_threshold:
        return "notify_admin"
    return "log_only"
\end{tcblisting}
